\begin{theorem}[Artin--Schreier phase]\label{mod:finitefield-artinschreier}
Let $k$ be a finite field with $\operatorname{char}k=2$, write
$\F_2=\mathbf Z/2\mathbf Z$, and let
\[
 \Tr=\Tr_{k/\F_2}:k\longrightarrow\F_2,
 \qquad
 \psi_0(x)=(-1)^{\Tr(x)}.
\]
In Lean these are \texttt{absoluteTraceTwo} and
\texttt{absoluteTraceChar}; the latter is obtained by pulling the exact
binary character $t\mapsto(-1)^t$ back along the trace.  It is nontrivial,
and the trace pairing is perfect.  Consequently every complex additive
character of $k$ is uniquely of the form $x\mapsto\psi_0(ax)$.

Both trace and $\psi_0$ are Frobenius invariant:
\[
 \Tr(x^2)=\Tr(x),
 \qquad
 \psi_0(x^2)=\psi_0(x).
\]
It follows, with the manuscript's exact normalization, that
\[
 \Tr(c+c^2)=0,
 \qquad
 \psi_0(c+c^2)=1
 \quad(c\in k).
\]

The file also proves the finite-field exactness needed later.  The
$\F_2$-linear Artin--Schreier map
\[
 \operatorname{AS}(c)=c+c^2
\]
has kernel $\F_2\cdot1=\{0,1\}$ and image the trace-zero hyperplane.  Thus
\[
 \Tr(a)=0\quad\Longleftrightarrow\quad
 \exists c\in k,\ a=c+c^2.
\]
If $\sigma(x)=x^2$ is Frobenius, then inverse Frobenius is its trace
adjoint:
\[
 \Tr(ax^2)=\Tr(\sigma^{-1}(a)x),
 \qquad
 \psi_0(ax^2)=\psi_0(\sigma^{-1}(a)x).
\]
Every trace-zero element is correspondingly of the form
$d+\sigma^{-1}(d)$.  Moreover, $\psi_0$ is the unique nontrivial additive
character annihilating every Artin--Schreier value.  For an arbitrary
nontrivial additive character $\psi$, Lean constructs a nonzero $c_0$ such
that
\[
 \psi(x^2)=\psi(c_0x),
 \qquad
 \psi(c_0^2)=\psi_0(1)=(-1)^{[k:\F_2]}.
\]
Finally, squaring is bijective on $k$; in particular every residue
coefficient has a unique square root.  This last finite-field fact is the
input used by the later reduced Artin--Schreier proof of odd break parity.
No local-field ramification or break-parity conclusion is assumed or proved
in this node.
\LeanFile{LanglandsFirstMainLemma/FiniteField/ArtinSchreier.lean}
\lean{LanglandsFirstMainLemma.artinSchreier_phase}
\leanok
\SourceText{Lemma lem:AS-phase-new and the finite-field perfectness and
Artin--Schreier input in the proof of Lemma lem:AS-break-parity; not its
local-field break-parity conclusion.}
\uses{mod:finitefield-quadraticphase}
\FormalizationContract
\end{theorem}
\begin{proof}
Frobenius is an $\F_2$-algebra automorphism of $k$, so invariance of
algebra trace gives the two Frobenius identities and hence the displayed
Artin--Schreier phase.  The kernel of $c\mapsto c+c^2$ consists of the two
roots of $c(c+1)$; rank--nullity and surjectivity of the finite-field trace
then identify its range with $\ker\Tr$.  Perfectness of the trace pairing
and finite Pontryagin duality give the additive-character parametrization
and rigidity statements.  Applying the same trace identity to a product
gives the inverse-Frobenius adjoint formula.  Finite-field perfectness gives
the unique square-root assertion.
\end{proof}
